% !PRINTED 101
\weberpage{101}

\begin{center}
{\Large\gesperrt{Third Section.}}

\vspace{1em}
{\large\gesperrt{The Roots of Algebraic Equations.}}
\end{center}

\begin{center}
{\large \S.\ 30.}

\vspace{0.5em}
{\large\gesperrt{Concept of Roots. Multiple Roots.}}
\end{center}

After, in the first two sections, the algebraic magnitudes had been considered more from the formal side, where the matter was identical transformations of letter-expressions in which the letters throughout could be conceived as symbols for variable magnitudes, the numerical magnitudes now come more into the foreground.

Here, in accordance with what was fixed in the Introduction, by numbers we understand real or imaginary magnitudes of the form $a+bi$, and for illustration we represent the real magnitudes by points of a straight line and the imaginary magnitudes by points of a plane. By the \gesperrt{absolute value} of an imaginary magnitude $a+bi$ we understand $\sqrt{a^2+b^2}$ and denote it, following Weierstrass, by
\[
 |a+bi|.
\]
Let now
\begin{equation}
 f(x)=a_0x^n+a_1x^{n-1}+a_2x^{n-2}+\cdots+a_n
\tag{1}
\end{equation}
be an integral function of $x$, in which the coefficients are arbitrary real or imaginary numbers, and the first, $a_0$, is assumed different from zero. If $\alpha$ is a number which, when put for $x$, makes the function $f(x)$ equal to zero, that is, which satisfies the condition $f(\alpha)=0$, then $\alpha$ is called a root of the equation
\[
 f(x)=0.
\]
